% Part: incompleteness
% Chapter: representability-in-q
% Section: prim-rec

\documentclass[../../../include/open-logic-section]{subfiles}

\begin{document}

\olfileid[fr]{inc}{req}{pri}
\olsection{Simulation de la récurrence primitive}

Nous pouvons maintenant montrer que la définition par récurrence
primitive se « simule » au moyen de la minimisation régulière et
de la fonction bêta. Supposons données les fonctions totales
$f(\vec x)$ et $g(\vec x,y,z)$. La fonction $h(\vec x,y)$ obtenue à partir de
$f$ et de~$g$ par récurrence primitive est définie par
\begin{align*}
h(\vec x, 0) & =  f(\vec x) \\
h(\vec x, y+1) & =  g(\vec x, y, h(\vec x, y)).
\end{align*}
Nous devons montrer que $h$ peut être définie à partir de $f$ et
de~$g$ en n'utilisant que la composition, la minimisation
régulière, les fonctions de base et les fonctions déjà obtenues
à partir de celles-ci par ces deux opérations, comme~$\beta$.

\begin{lem}
\ollabel{lem:prim-rec}
Si $h$ se définit à partir des fonctions totales $f$ et $g$ par
récurrence primitive, elle se définit aussi à partir de $f$, de $g$ et des
fonctions $\Zero$, $\Succ$, $\Proj{n}{i}$, $\Add$, $\Mult$ et
$\Char{=}$, au moyen de la composition et de la minimisation
régulière.
\end{lem}

\begin{proof}
Définissons d'abord une fonction auxiliaire $\hat h(\vec x,y)$
qui renvoie le plus petit nombre~$d$ codant une suite telle que
\begin{enumerate}
\item $(d)_0=f(\vec x)$, et
\item pour chaque $i<y$, $(d)_{i+1}=g(\vec x,i,(d)_i)$,
\end{enumerate}
où $(d)_i$ abrège maintenant $\beta(d,i)$. Autrement dit,
$\hat h$ renvoie le code d'une suite commençant par
$\tuple{h(\vec x,0),h(\vec x,1),\dots,h(\vec x,y)}$.
Nous pouvons écrire
\[
\hat h(\vec x, y) = \umin{d}{(\beta(d,0) = f(\vec x) \land \bforall{i <
  y}{\bigl(\beta(d,i+1) = g(\vec x,i,\beta(d,i))\bigr)})}.
\]
Aucune récurrence primitive n'est nécessaire ici : il suffit
d'une minimisation. Le lemme de la fonction bêta
\olref[bet]{lem:beta} garantit que la recherche est régulière,
car toute suite finie des valeurs prescrites possède un code.

Nous avons alors
\[
h(\vec x, y) = \beta(\hat h(\vec x, y), y),
\]
de sorte que $h$ se définit à partir de $f$, de $g$ et des
fonctions de base par composition et minimisation régulière.
\end{proof}

\begin{explain}
La condition testée par la minimisation fixe les valeurs
décodées aux positions $0,\dots,y$. Une récurrence sur ces
positions montre qu'elles valent $h(\vec x,0),\dots,h(\vec x,y)$
pour tout code qui passe le test. Plusieurs codes peuvent
convenir. La minimisation en choisit un et le décodage à la
position~$y$ donne la même valeur pour chacun d'eux.
\end{explain}

\end{document}
